\section{Audience, scope, and status terminology}
\label{sec:orientation}

\subsection{What this guide makes easier to read}

The existing \sourcepath{latex/} documents in this workspace are weekly
research and development records. This guide reorganizes the material so that
mathematicians and cryptographers can directly answer three questions.

\begin{enumerate}
  \item Which mathematical equality in the paper becomes which byte invariant in the implementation?
  \item What is the strongest theorem actually checked by EasyCrypt so far, and what are its premises?
  \item What else is needed for that theorem to become complete functional correctness or an EUF-CMA theorem?
\end{enumerate}

Actual EasyCrypt declarations are the primary evidence in this guide.
\sourcepath{CLAIM_LEDGER.md} and \sourcepath{manifests/paper-artifacts.md}
define the frozen state; post-freeze reports take precedence for subsequent
results. Extraction manifests and weekly reports follow. The opening portions
of weekly documents are historical snapshots, so the following materials take
precedence when assessing the current state.

\begin{itemize}[leftmargin=1.5em,itemsep=0.1em,topsep=0.3em]
  \raggedright
  \item \sourcepath{manifests/paper-artifacts.md},
        \sourcepath{POST_FREEZE_KG_RQ_}\allowbreak
        \sourcepath{HAETAE_REPORT.md},
        \sourcepath{POST_FREEZE_KG_ACTUAL_}\allowbreak
        \sourcepath{AVEC_QJ_REPORT.md};
  \item \sourcepath{WEEK16_KG_NTT_MUL_}\allowbreak\sourcepath{REPORT.md},
        the Week~16 KG-NTT audit note \sourcepath{latex/sections/}\allowbreak
        \sourcepath{33-week16-kg-ntt-mul-stop.tex};
  \item \sourcepath{WEEK16_SIGN_REPORT.md},
        \sourcepath{WEEK16_VERIFY_REPORT.md},
        \sourcepath{logs/verify-all-week16-}\allowbreak
        \sourcepath{verify-matrix-crt.log};
  \item \sourcepath{WEEK16_KG_REPORT.md},
        \sourcepath{WEEK16_MINCORE_PLAN.md},
        \sourcepath{logs/verify-all-week16-}\allowbreak
        \sourcepath{kg-first-task.log};
  \item \sourcepath{WEEK15_REPORT.md},
        \sourcepath{RANS_ACTUAL_SUCCESS_}\allowbreak\sourcepath{WITNESS.md},
        the Week~15 research note \sourcepath{latex/sections/}\allowbreak
        \sourcepath{31-week15-rans-actual-}\allowbreak
        \sourcepath{success-witness.tex};
  \item \sourcepath{WEEK14_REPORT.md}, the Week~14 research note
        \sourcepath{latex/sections/}\allowbreak
        \sourcepath{30-week14-hbz-full-wrapper-}\allowbreak
        \sourcepath{composition.tex},
        \sourcepath{HBZ_FULL_WRAPPER_}\allowbreak\sourcepath{COMPOSITION.md};
  \item \sourcepath{WEEK13_REPORT.md}, the Week~13 research note
        \sourcepath{latex/sections/}\allowbreak
        \sourcepath{29-week13-rans-core-}\allowbreak\sourcepath{composition.tex},
        the independent verification record \sourcepath{logs/week13-}\allowbreak
        \sourcepath{independent-verifier.md};
  \item \sourcepath{PRODUCT_REPLAY_}\allowbreak\sourcepath{DECISION.md},
        \sourcepath{SIGNATURE_CODEC_}\allowbreak\sourcepath{OBLIGATIONS.md},
        \sourcepath{MODE2_HBZ_CODEC_}\allowbreak\sourcepath{OBLIGATIONS.md};
  \item \sourcepath{RANS_PROOF_}\allowbreak\sourcepath{DESIGN.md},
        \sourcepath{RANS_BYTE_STACK_}\allowbreak\sourcepath{INVARIANT.md},
        \sourcepath{RANS_ENCODER_}\allowbreak\sourcepath{INVARIANT.md},
        \sourcepath{RANS_DECODER_}\allowbreak\sourcepath{INVARIANT.md}, and operational documents.
\end{itemize}

Line numbers change during editing, so the guide uses file paths and
declaration names as stable identifiers.

\subsection{Terminology conventions}

The original Korean edition used the following conventions to help
mathematicians and cryptographers match implementation terminology to the
source. This English translation uses the corresponding English terms.

\begin{enumerate}
  \item On first use, the Korean term preceded its English original in parentheses.
        Examples include refinement, invariant, trace, frame condition, prefix, and suffix.
  \item Repeated uses in the same context preferred Korean, while sentences
        directly compared with code could repeat the previously introduced English term.
  \item EasyCrypt theorem and procedure names, claim identifiers, filenames,
        and code snippets were left unchanged for exact lookup. Standard
        abbreviations such as \identifier{rANS}, \identifier{HBZ},
        \identifier{W32}, \identifier{W64}, and \identifier{EUF-CMA} are also retained.
  \item The terms ``actual'' for an actual generated procedure, ``canonical,''
        and ``success-conditioned'' were introduced with their Korean counterparts.
\end{enumerate}

\subsection{Status terms}

\begin{description}[style=nextline,leftmargin=2.7cm]
  \item[\statusproved]
    The EasyCrypt theorem compiles without cached proofs under the stated
    premises. This does not mean that every premise of the parent goal is discharged.
  \item[\statuspartial]
    Genuine subtheorems needed by the parent goal exist, but unresolved
    premises or call boundaries remain for the complete parent goal.
  \item[\statusspecified]
    The goal's type and logical formula compile, but no theorem proving it is yet available.
  \item[\statusblocked]
    A required extraction, memory semantics, probability distribution,
    numerical-error bound, or cryptographic reduction is concretely missing,
    preventing the next composition.
  \item[\statusdeferred]
    The gap is precisely identified but is not closed under the currently
    permitted proof rules and scope, so it is frozen as an explicit paper
    boundary. This does not mean that it is proved or unnecessary.
\end{description}

\begin{boundary}[Meaning of \texttt{PROVED}]
In this guide, \statusproved always means that the stated theorem has been
verified for executions satisfying its displayed premises. In particular,
\statusproved for a Hoare partial-correctness theorem does not imply
termination, losslessness, or equality of output distributions.
\end{boundary}

\subsection{Snapshot and targets}

\begin{table}[ht]
\centering
\caption{Verification boundaries fixed by this guide}
\label{tab:snapshot}
\begin{tabularx}{\textwidth}{L{0.27\textwidth}Y}
\toprule
Item & Fixed target \\
\midrule
Paper specification & \sourcepath{CryptoLab_HAETAE.pdf}, 46 pages, created 2026-02-04,
SHA-256 \\
& \identifier{80d28b3af9af2a27dba0eb4746f9b575}\newline
\identifier{5e05cb9a86bb6596f45fda06c605f3f6}
\cite{haetae-spec} \\
Parameter set & HAETAE-2 only: \(n=256\), \((k,\ell)=(2,4)\),
\(q=64513\), \(\eta=1\), \(\tau=58\) \\
Implementation & Pinned \sourcepath{haetae-ref-jasmin/} sources and
\identifier{jasmin2ec} outputs \cite{haetae-jasmin} \\
Frozen authored proofs & 82 manifested EasyCrypt targets under
\sourcepath{haetae-topdown-easycrypt/easycrypt/}; both the pre-freeze 82/82 and
current frozen-suite 82/82 results are preserved \\
Frozen paper decision & \identifier{PAPER-FROZEN}; internal stopping points are
\identifier{STOP-KG-NTT}, \identifier{STOP-SIGN-CHAL-MODE2},
\identifier{STOP-VERIFY-MATRIX-CRT} \\
Latest post-freeze decision & \identifier{STOP-KG-AVEC-QJ};
\identifier{OBL-KG-SECURITY-EXPANDVECA-SEMANTICS} is the first repair obligation.
Pure NTT/Rq and actual-\(a\)/paper-\(qj\) auxiliary theorems are verified
separately outside the frozen manifest \\
Reused proofs & Import-only adapters retaining the original premises of pinned
KeyGen, NTT, sampler, singularity, and FFT results \cite{haetae-easycrypt} \\
Excluded scope & Other parameter sets, physical side channels, faults, RNG
quality, and memory safety and SCT that do not follow automatically from EUF-CMA \\
\bottomrule
\end{tabularx}
\end{table}

\begin{boundary}[The frozen 82 targets and post-freeze auxiliary evidence]
After target 77, \sourcepath{Mode2KeygenNttMulBridge.ec}, target 78,
\sourcepath{Mode2SignAcceptedCore.ec}, and target 82,
\sourcepath{Mode2VerifyCoreSequence.ec}, freshly compile the actual Sign/Verify
control boundaries. The completed 78-target run before the Verify changes is
preserved in \sourcepath{logs/verify-all-before-week16-verify.log}. Preserved
\sourcepath{logs/verify-all-before-paper-freeze.log}, with final line
\texttt{RESULT PASS}\newline
\texttt{authored-targets=82 cache=-no-eco} and SHA-256
\hashsplit{4cd64e5a656be82710bca1410c4d194}{03a3c661d6b91b0319a0ea8f7c91646da},
is completed evidence immediately before the freeze. The
\sourcepath{logs/verify-all-summary.txt} tracked at baseline commit
\identifier{e87bcc3} includes the paper-freeze audit and has the same terminal
result with SHA-256
\hashsplit{08ef9639dc73d56dba42d02999d07897}{d29bc8e60aa100a648e9437fd64387ad}.
The three \sourcepath{post-freeze/KgActualAvecQj*.ec} targets are not included
in these 82 and were each freshly compiled separately. \statusproved for a
Hoare theorem still does not imply actual procedure termination or probability-1 execution.
\end{boundary}

\begin{takeaway}{What this guide does not claim}
The work does not yet claim complete KeyGen distribution equality, full
Sign/Verify functional correctness, equality of public-context transcripts,
direct \(\mu\) equality for accepted raw-API traces, the complete 1474-byte
signature codec, actual rANS success reachability for general canonical inputs,
encoder/decoder termination, unconditional round trips of the actual complete
HBZ wrappers on every canonical input, the (h) codec, identification with
full-NTT/paper multiplication within the frozen claim matrix, identification of
the security model's ExpandVecA with actual \identifier{avec}, the paper's
KeyGen equation, MINCORE Sign/Verify composition, encoding zero loss, or
end-to-end EUF-CMA security of the Jasmin implementation. The success
postcondition for fixed all-zero input is proved, but termination/losslessness
of that execution is separate.
\end{takeaway}
